% 本文件由 LaTeX 排版 Agent 根据有效 translation.json 自动生成。
% author=Apocrypha; work_id=0849; title=玛利亚诞生福音

\chapter{玛利亚诞生福音}

% structure: p0001 source=h1 target=omitted reason=page-title-already-rendered-by-parent

% structure: p0002 source=h2 target=section reason=relative-heading-rank; base=h2

\section{第一章}

蒙福而荣耀的卒世童贞玛利亚，出自达味的王族和家系，生于纳匝肋城，在耶路撒冷上主的圣殿中受教养。她的父亲名叫约阿希姆，母亲名叫亚纳。父亲的家系来自加里肋亚和纳匝肋城，母亲的家系则来自白冷。他们的生活在上主面前纯全正直，在人面前无可指摘且虔诚。因为他们将全部财产分为三份：一份用于圣殿和圣殿的仆役；另一份分给旅客和穷人；第三份留给自己和家庭的所需。这样，他们爱天主，善待他人，在自家房屋中度着贞洁的婚姻生活约二十年，没有子女。然而他们誓愿，倘若上主赐予他们后裔，就将孩子奉献于上主侍奉；因此他们每年在节期时都前往上主的圣殿。\par

% structure: p0004 source=h2 target=section reason=relative-heading-rank; base=h2

\section{第二章}

到了献殿节临近的时候，约阿希姆便与本支派的几个人一同上耶路撒冷去。那时，依撒加尔在那里任大司祭。依撒加尔看见约阿希姆与其他同乡一同献祭，就轻视他，拒绝他的礼物，问他这没有后代的人怎敢站在那些有子女的人中间；并说他的礼物绝不能为天主所悦纳，因为天主认为他不配拥有后裔：因为经上记载说：\Quote{凡在以色列中没有生男或生女的人，都是受诅咒的。}因此他说，他应当先藉生育子女摆脱这诅咒，然后才可带着祭品来到上主面前。约阿希姆因这当面羞辱而满面羞愧，便退到牧人那里，他们正在牧场牧放羊群；他也不愿回家，恐怕本支派的人也会同样羞辱他，因为他们当时在场，并听到了司祭的话。\par

% structure: p0006 source=h2 target=section reason=relative-heading-rank; base=h2

\section{第三章}

他在那里停留了一段时日。某一天，当他独自一人时，上主的一位天使在极大的光明中站在他身边。约阿希姆因天使的出现而惊慌，显现的天使便平息他的恐惧说：\Quote{约阿希姆，不要害怕，也不要因我的显现而惊慌；因为我是上主的天使，奉祂差遣到你这里来，告诉你：你的祈祷已蒙垂听，你的施舍已上升到祂面前。\BibleRef{宗 10:4} 因为祂看见了你的羞辱，听见了那不生育的责难，这是不公正地加给你的。因为天主是罪的报复者，而非本性的报复者；因此，当祂封闭任何人的子宫时，祂这样做是为了以后奇迹般地再次开启，使所生的被承认不是出于情欲，而是出于天主的恩赐。难道你们民族的第一个母亲撒辣不是直到八十岁还不生育吗？然而，她在极其高龄时生下了依撒格，万民的祝福许诺重新赐给了他。辣黑耳也是如此，她深得上主宠爱，为圣雅各伯所钟爱，却长期不育；但她生下了若瑟，他不仅是埃及的主宰，也是许多濒临饿死之人的拯救者。在民长中，有谁比三松更强壮，或比撒慕尔更圣洁呢？然而，两者的母亲都是不育的。因此，如果我的话语的合理性不能说服你，就当相信事实：晚年怀孕和不育妇女的生育，通常伴随着某种奇妙之事。因此，你的妻子亚纳将为你生一个女儿，你要给她起名叫玛利亚；她将如你所许愿的，自幼奉献给上主，并从母胎中就充满圣神。她不会吃或喝任何不洁之物，也不会在外面的民众中间度过一生，而是在上主的圣殿中，使人无法说她或甚至怀疑她有什么恶事。因此，当她长大后，正如她自己要奇迹般地由不育妇人所生，同样，她——一位童贞女——要以无可比拟的方式生下至高者的儿子，他要被称为耶稣，按照祂名字的含义，祂是万民的救主。这就是我所宣告之事的记号：当你来到耶路撒冷的金门时，你会在那里遇见你的妻子亚纳，她因你迟归而焦虑，但见到你时会满心欢喜。}天使说完这话，就离开了他。\par

% structure: p0008 source=h2 target=section reason=relative-heading-rank; base=h2

\section{第四章}

此后，天使向他妻子亚纳显现说：\Quote{亚纳，不要害怕，也不要以为你所看见的是幻象。因为我是那位将你的祈祷和施舍呈于天主前的天使；现在我奉差遣到你这里，向你宣告：你将生一个女儿，她要叫玛利亚，她要在众女人中蒙受祝福。她自出生就充满上主的恩宠，要在父亲家中住三年，直到断奶。此后，她被交付于上主的侍奉，在圣殿中不得离开，直到达至明辨的年龄。总之，她在那里日夜以斋戒和祈祷侍奉天主，不沾任何不洁之物；她永不认识男人，而是独一无二、无玷、无染、不与男人交往，她作为童贞女要生一个儿子；她作为祂的婢女，要生下主——无论是在恩宠、名字还是事工上，祂都是世界的救主。因此，起身，上耶路撒冷去；当你来到那因镀金而称为金门的门时，作为记号，你会在那里遇见你的丈夫，你曾为他的安全焦虑。当这些事如此发生时，你就知道我所宣告的必定成就。}\par

% structure: p0010 source=h2 target=section reason=relative-heading-rank; base=h2

\section{第五章}

因此，正如天使所吩咐的，他们二人从所在之地出发，上到耶路撒冷；当他们来到天使预言指示的地方时，便在那里相遇了。于是，他们因彼此相见而欢喜，且因对所应许的后裔确信不疑，便向那举扬谦卑者的主献上当得的感谢。如此，朝拜了主之后，他们返回家中，怀着确信与喜乐等候天主的应许。安娜于是怀孕，生了一个女儿；按照天使的命令，她的父母给她起名叫玛利亚。\par

% structure: p0012 source=h2 target=section reason=relative-heading-rank; base=h2

\section{第六章。}

当三年之期届满，离乳之时满足，他们便带着祭品将这童贞女带到主的圣殿。圣殿周围按着十五首《阶梯圣咏》有十五级台阶上升；因圣殿建在山上，外边的全燔祭台须经台阶方能到达。于是，她的父母将小女孩，即蒙福的童贞玛利亚，放在其中一级上。当他们脱下旅途所穿的衣服，照例换上更整洁干净的衣服时，主的童贞女一级一级地登上所有台阶，无人搀扶或托举，以至于就这一点而言，你会以为她已成年。因为主在祂的童贞女幼年时已行了一件大事，并以此奇迹的征兆预示了她将来是何等伟大。因此，按律法的惯例献了祭，完成了他们的誓愿，他们便将童贞女留在圣殿的围墙内，与其他童贞女一同受教育，而他们自己则返回家中。\par

% structure: p0014 source=h2 target=section reason=relative-heading-rank; base=h2

\section{第七章。}

但主的童贞女在年龄和德行上渐长；虽然按圣咏作者的话，她的父母遗弃了她，但主却收留了她。因为她每日有天使拜访，每日享受神视，这保守她远离一切邪恶，使她充满一切美善。如此，她到了十四岁；不仅恶人无法指控她任何可责备的事，凡认识她的生活和言行的善人，都认为她堪受钦佩。于是大司祭公开宣告，那些公开住在圣殿中、已达此年龄的童贞女，应按民族习俗和年龄成熟回家结婚。其他人都欣然服从这命令；唯独玛利亚，主的童贞女，回答说她不能这样做，说她的父母已将她献于主的事奉，而且她自己也向主许了童贞的誓愿，绝不会因与任何男人交往而违反此誓。大司祭心中十分困惑，一方面他认为不应违背圣经而破坏誓愿，经上说：\Quote{你许了愿，就要偿还}；另一方面他也不敢引入民族所不知的习俗。于是他下令，在即将来临的庆节上，耶路撒冷及其附近的所有首领都应出席，以便凭他们的建议知道在此疑难之事上该如何行。这事发生后，他们一致决定应就此事求问主。当他们全体俯首祈祷时，大司祭照常去求问天主。没有等待多久，就在众人听闻中，从神谕所和赎罪盖那里发出声音说：按照依撒意亚的预言，应寻找一个人，童贞女应托付给他并与他订婚。因为依撒意亚明说：\BibleQuote{从叶瑟的根上将发出一条枝条，从他的根上将升起一朵花；上主的神将停在他身上，智慧和聪敏的神，超见和刚毅的神，明达和虔敬的神；他将充满敬畏上主的神。} \BibleRef{依 11:1} 根据这预言，他因此预言所有达味家族和亲属中未婚且适合结婚的男子，都应带他们的枝条到祭台前；谁的枝条带来后开了花，并且枝条末端有上主的神以鸽子形状停驻，那人就是童贞女应托付给他并与之订婚的人。\par

% structure: p0016 source=h2 target=section reason=relative-heading-rank; base=h2

\section{第八章。}

在众人中有一位若瑟，属于达味的家族和亲属，年事已高；当所有人按命令带来枝条时，唯独他留下了自己的。因此，当没有出现与神声相符之事时，大司祭认为有必要第二次求问天主；天主回答说，在那些被指定的人中，唯独那位应与之订婚的人没有带他的枝条。于是若瑟被找到了。因为他带来了他的枝条，鸽子从天降下，落在枝条顶端，众人清楚看出他就是童贞女应与之订婚的人。因此，举行了通常的订婚仪式后，他回到白冷城，整理家务，预备婚礼所需之物。但玛利亚，主的童贞女，与另外七位和她同龄、同时离乳的童贞女（她是从司祭那里领受的）一起，返回了加里肋亚她父母的家中。\par

% structure: p0018 source=h2 target=section reason=relative-heading-rank; base=h2

\section{第九章。}

在那些日子里，就是她初到加里肋亚的时候，天主派遣天使加俾额尔到她那里，向她报告主受孕的事，并解释受孕的方式和次序。于是他进去，使她所处的房间充满了极大的光明；并极其谦恭地向她请安，说：\Quote{万福，玛利亚！极蒙主宠爱的童贞女，充满恩宠的童贞女，主与你同在；你在一切妇女中受赞颂，你在至今所生的一切人中受赞颂。}\BibleRef{路 1:26}童贞女早已熟悉天使的面容，也不惊异于天上的光明，所以既不被天使的显现吓倒，也不因光明的伟大而惊愕，只是因他的话而困惑；她开始思忖，这样不寻常的请安是什么性质，预示什么，或有什么目的。天使受天主启示，接过这心思，说：\Quote{玛利亚，不要害怕，好像这请安下隐藏着什么违反你贞洁的事。因为你在选择贞洁时，已在主前蒙恩；因此你，一位童贞女，将无玷受孕，生一个儿子。他将伟大，因为他将统治从海到海，从大河直到地极；他将被称为至高者之子，因为他在地上卑微地出生，在天上高升地为王；主天主必把他祖先达味的宝座赐给他，他要为王统治雅各伯家直到永远，他的王国无穷无尽；\BibleRef{路 1:32}因为他是万王之王，万主之主，\BibleRef{默 19:16}他的宝座从永远到永远。}童贞女没有怀疑天使的话；但愿知道其方式，就回答说：\Quote{这事怎能成就？因为按我的誓愿，我不认识男人，怎能不加入男人的种子而生呢？}天使回答说：\Quote{玛利亚，不要以为你将如同人类那样怀孕；因为与男人没有任何交合，你，一位童贞女，将怀孕；你，一位童贞女，将生产；你，一位童贞女，将哺乳；因为圣神要临于你，至高者的能力要庇荫你，\BibleRef{路 1:35}毫无情欲的热火；因此你所生的那一位，将是唯一圣洁的，因为只有他，无玷受孕和出生，才被称为天主之子。}于是玛利亚伸出双手，举目向天，说：\Quote{请看，主的婢女，因为我不配称为主母；愿照你的话成就于我吧。}\par

将会很长，也许对某些人来说甚至是乏味的，如果我们在本小作品中插入所有我们在主诞生前后读到的事；因此，省略那些在\WorkTitle{福音}中更充分地记载的事，就让我们来谈那些被认为不太值得叙述的事。\par

% structure: p0021 source=h2 target=section reason=relative-heading-rank; base=h2

\section{第十章}

因此，若瑟从犹太来到加里肋亚，要娶与他订婚的童贞女为妻；因为已经过了三个月，自她与他订婚以来，第四个月刚开始。同时，从她的身量明显看出她怀孕了，她也不能向若瑟隐瞒此事。因为他与她订婚了，可以更自由地来见她，更亲切地与她谈话，他就发现她有了身孕。于是他开始非常犹豫和困惑，因为他不知道怎么做最好。因为他是义人，不愿意公开羞辱她；又因他是虔诚人，不愿因怀疑行淫而损害她的美名。因此他决定私下解除婚约，暗暗休了她。正当他思虑这些事时，看，主的天使在梦中显现给他，说：\Quote{若瑟，达味之子，不要害怕，就是说，不要对这童贞女有任何行淫的怀疑，也不要以为她有什么恶事；不要害怕娶她为妻，因为她腹中所孕的、现在使你忧愁的，不是男人的作为，而是圣神的作为。因为她在所有童贞女中，唯独要生下天主之子，你要给他起名叫耶稣，即救世主，因为他要把自己的百姓从罪恶中救出来。}因此，若瑟按照天使的吩咐，娶了这位童贞女为妻；但没有认识她，却照顾她，并保守她的贞洁。\BibleRef{玛 1:18}现在她怀孕的第九个月近了，若瑟便带着他的妻子和所需之物，往他本城白冷去。他们在那里的时候，她分娩的日期满了，她就生了头胎儿子，正如圣史们所显示的，就是我们的主耶稣基督，他与圣父、圣子、圣神，永生永王，是天主，从永远到永远。\par

\SourceRecord{Translated by Alexander Walker.；Ante-Nicene Fathers；Vol. 8.；Edited by Alexander Roberts, James Donaldson, and A. Cleveland Coxe.；Buffalo, NY: Christian Literature Publishing Co.,；1886.；<http://www.newadvent.org/fathers/0849.htm>.}
